\abstract{\textbf{Motivation:} Macrocyclic peptides are promising therapeutic candidates for intracellular targets, but their design requires simultaneous control over non-natural monomer chemistry, ring topology, membrane permeability, and target binding. Existing SMILES- or HELM-string generative models either operate in long atom-level sequence spaces or treat monomers as symbolic tokens with limited chemical grounding.\\
\textbf{Results:} We introduce \mytool, an AutoRegressive Latent Diffusion(ALD) framework for \textit{de novo} macrocyclic peptide generation. The model represents HELM monomers with structured Uni-Mol embeddings, generates each residue through context-conditioned diffusion in chemical latent space, predicts R-group-aware ring closures during autoregressive generation, and aligns the denoiser to permeability--affinity rewards using winner-protected diffusion-adapted preference optimization. In silico experiments demonstrate \mytool's generation quality and reward-optimization performance against representative peptide generation baselines.\\
\textbf{Availability:} Code and data will be made available in the project repository.\\
\textbf{Contact:} \href{email@example.com}{email@example.com}\\
\textbf{Supplementary information:} Supplementary data are available at \textit{Bioinformatics} online.}
